
Before continuing the proof of \SGARefLink{sga3:xvii:subsection:4:1}{4.1}, note that the equivalence
(i)\(\Leftrightarrow\)(ii), together with \SGARefLink{sga3:xvii:lemma:2:3:bis}{2.3 bis}, gives the
following corollary.

\SGAStatement{Corollary 4.1.3}{If \(k\) is a perfect field, a
smooth, connected \(k\)-algebraic group is unipotent if and only if
it has a composition series whose successive quotients are
isomorphic to \(\mathbb G_a\).}
\SGARefTarget{sga3:xvii:corollary:4:1:3}

\SGAProof{Continuation of the proof of \SGARefLink{sga3:xvii:subsection:4:1}{4.1}}
The implication (i)\(\Rightarrow\)(iv) follows from \SGARefLink{sga3:xiv:2.4:item:i}{2.4(i)}.

For (iv)\(\Rightarrow\)(v), by \SGARefLink{sga3:xiv:4.1}{Exposé XIV, 4.1}, the group \(G\)
has a maximal torus \(T\) defined over \(k\).  If \(r=\dim T\),
then
\[
  ({}_nT)_{\overline{k}}\simeq(\mu_n)^r.
\]
It follows that \(r=0\).

The implication (v)\(\Rightarrow\)(i) follows from the observation
made in the proof of \SGARefLink{sga3:xvii:corollary:3:8}{3.8}.

For (i)\(\Rightarrow\)(vi), by \SGARefLink{sga3:xvii:corollary:3:4}{3.4} the group \(G\) is in fact
contained in an algebraic subgroup of \(\operatorname{GL}(n)\)
isomorphic to \(\operatorname{Trigstr}(n)\).  Hence
\(\mathfrak g\) consists of nilpotent endomorphisms.

For (vi)\(\Rightarrow\)(vii), the Lie algebra \(\mathfrak g\) is
nilpotent by Engel's theorem (Bourbaki, \emph{Groupes et algèbres
de Lie}, Chapter I, \S4, Corollary 3).

For (vii)\(\Rightarrow\)(v), let \(T\) be a maximal subtorus of
\(G\) (\SGARefLink{sga3:xiv:1.1}{Exposé XIV, 1.1}), and let \(\mathfrak t\) be its Lie
algebra.  The embedding of \(G\) in \(\operatorname{GL}(n)\)
defines a representation of \(T\), which is necessarily
semisimple; this can be seen after passing to an algebraic closure
\(\overline{k}\) of \(k\) and applying \SGARefLink{sga3:i:proposition:4-7-3}{Exposé I, 4.7.3}.  Thus, if
\(X\in\mathfrak t\), then \(X\) is a semisimple endomorphism in
\(\mathfrak{gl}(n)\).  It follows at once that
\[
  \operatorname{ad}X:\ Y\longmapsto[X,Y]
\]
is a semisimple endomorphism of \(\mathfrak{gl}(n)\), and hence of
\(\mathfrak g\).  This endomorphism is also nilpotent, since
\(\mathfrak g\) is nilpotent.  Therefore
\(\operatorname{ad}X=0\), so \(X\) is central.  Consequently
\(\mathfrak t\) is central and consists of semisimple
endomorphisms; by hypothesis it is zero.  A fortiori, \(T\) is the
trivial group.

\SGAStatement{Remark 4.1.4}{%
\begin{enumerate}
\item[(a)] \SGARefTarget{sga3:xvii:remark:4:1:4:item:a}
  We shall give below (\SGARefLink{sga3:xvii:proposition:4-3-1}{4.3.1}) an infinitesimal characterization of
  unipotent groups in characteristic \(p>0\) that is independent of
  an embedding in \(\operatorname{GL}(n)\).
\item[(b)] \SGARefTarget{sga3:xvii:remark:4:1:4:item:b}
  When \(k\) is perfect, conditions (ii) and (iii) of \SGARefLink{sga3:xvii:proposition:4-1-1}{4.1.1}
  simplify by virtue of the following lemma.
\end{enumerate}}
\SGARefTarget{sga3:xvii:remark:4:1:4}

\SGAStatement{Lemma 4.1.5}{If \(k\) is a perfect field, every
\(k\)-algebraic group \(G\) that is a form of
\((\mathbb G_a)^r\) is isomorphic to \((\mathbb G_a)^r\).}
\SGARefTarget{sga3:xvii:lemma:4:1:5}
